\section{The opening IOP}\label{sec:iopp}

The PIOP hands this section exactly three open claims:
\[
  E_{\rm eq}:\ipq{c\,\eq(\cdot,r_h)}{\pi_q(\hh)}=\tau,
  \qquad
  H:\pitwo(\hh)=M\pitwo(\ff^{\circ}),
  \qquad
  C:\orc=\cw{\ff}.
\]
The virtual witness $\hh$ is already a bit vector, so there is no
bit-decomposition step for SHA-256.  Crucially, the PIOP has already changed the
arbitrary coefficient vector $\vv=\bar A^{\!\top}\eq(\cdot,r_{\rm con})$ into
the evaluation-shaped vector $c\,\eq(\cdot,r_h)$.  The opening may now reshape
that vector without assuming a tensor factorisation of $\vv$.

\subsection{Axes and fields}

The committed vector $\ff$ already has the commitment-time matrix view
$F\defeq\operatorname{reshape}(\ff)$ from \cref{sec:commit}.  For the virtual
vector $\hh$, pad to a rectangular power-of-two layout and set
\[
  \begin{gathered}
    \Rows=\bits^{\drow},\qquad \Cols=\bits^{\dcol},\\
    \ellrow=2^{\drow},\qquad
    \ellcol=2^{\dcol},\qquad
    \ellrow\ellcol=2^{d_h}\geq|\hh|,
    \qquad d_h=\drow+\dcol.
  \end{gathered}
\]
The implementation calls these dimensions $\ell_1$ and $\ell_2$; this section
uses the semantic names $\ellrow$ and $\ellcol$.  Reshape the virtual vector as
\[
  H[\rowidx,\colidx]
  =\hh[\colidx\ellrow+\rowidx].                 \tag{4.1}\label{eq:reshape}
\]
The row/fold index is the fastest-moving coordinate, so each
$H[*,\colidx]$ is contiguous and a $128$-bit pack never crosses a column
boundary.  Split the sumcheck point according to exactly these coordinates,
\[
  r_h=(r_{\rm row},r_{\rm col}),
  \qquad
  r_{\rm row}\in\Fq^{\drow},\quad
  r_{\rm col}\in\Fq^{\dcol}.
\]
The equality polynomial factors over a concatenated point:
\[
  \eq\bigl((\rowidx,\colidx),r_h\bigr)
  =\eq(\rowidx,r_{\rm row})\eq(\colidx,r_{\rm col}).
\]
Define the two public weight vectors
\[
  v_{\rm row}[\rowidx]\defeq\eq(\rowidx,r_{\rm row}),
  \qquad
  v_{\rm col}[\colidx]\defeq c\,\eq(\colidx,r_{\rm col}).
\]
Then the PIOP boundary is exactly
\[
  E_{\rm eq}:
  \sum_{\rowidx\in\Rows}\sum_{\colidx\in\Cols}
     v_{\rm row}[\rowidx]v_{\rm col}[\colidx]\,
     \pi_q(H[\rowidx,\colidx])=\tau
  \quad\text{in }\Fq .                         \tag{4.2}
\]
This separation is an identity of the equality polynomial.  It is the output
of \cref{step:eqshape}, not a property assumed of the original $\vv$.
\begin{center}
\small
\(
  \underbrace{
    \operatorname{reshape}\!\left(
      \bar A^{\!\top}\eq(\cdot,r_{\rm con})
    \right)
  }_{\text{usually rank \(>1\)}}
  \quad
  \xRightarrow{\text{\cref{step:eqshape}: sumcheck}}
  \quad
  \underbrace{
    \bigl[\eq(\rowidx,r_{\rm row})\bigr]_{\rowidx}\,
    \bigl[c\,\eq(\colidx,r_{\rm col})\bigr]_{\colidx}^{\!\top}
  }_{\text{rank one for every split}} .
\)
\end{center}
The names $\rowidx$ and $\colidx$ always denote the matrix axes in
\eqref{eq:reshape}.  Inner products instead advertise their field:
$\ipq{\cdot}{\cdot}$, $\ipZ{\cdot}{\cdot}$, or
$\ipC{\cdot}{\cdot}$.

The field path is
\[
  \Fq\ \xrightarrow{\text{canonical lift}}\ \Z
  \ \xrightarrow{a\mapsto\gen^a}\ \comfield^{\times}
  \ \xrightarrow{\text{linear opening}}\ \comfield.
\]

% =========================================================================
\begin{step}{Reshape and send the column folds as integers}{%
  \clgdone{$E_{\rm eq}$}%
  \clgnew{T}{$\mu[\colidx]=\ipZ{\gam}{H[*,\colidx]}$}%
  \clgnote{T}{$\rowidx$ is summed; one integer remains per $\colidx$}%
}\label{step:tensor}
The core opening interface needs public weights that separate across the two
axes.  The preceding sumcheck guarantees exactly that form.  Canonically lift
its row coefficients from $\Fq$ to integers,
\[
  \gam[\rowidx]
  \defeq\pi_{\mathrm{can}}^{-1}
     \bigl(v_{\mathrm{row}}[\rowidx]\bigr)\in[0,q-1].
\]
For every fixed column, the prover sends the integer
\[
  \boxed{
  \mu[\colidx]
  \defeq
  \ipZ{\gam}{H[*,\colidx]}
  =\sum_{\rowidx\in\Rows}
       \gam[\rowidx]H[\rowidx,\colidx] .}       \tag{4.3}\label{eq:mucol}
\]
This dot product runs \emph{down the row axis} of one fixed column, and it runs
over $\Z$.  The verifier checks the original field claim by reducing the sent
integers modulo $q$:
\[
  \boxed{\ipq{v_{\mathrm{col}}}{\pi_q(\bm\mu)}=\tau.}  \tag{4.4}\label{eq:murecombine}
\]
Equations \eqref{eq:mucol} and \eqref{eq:murecombine} are exactly
$E_{\rm eq}$ written column by column.

\begin{center}
\begin{tikzpicture}[
  font=\scriptsize,
  cell/.style={draw=black!38, minimum width=0.48cm, minimum height=0.38cm,
               inner sep=0pt, line width=0.3pt},
  arr/.style={-{Latex[length=1.8mm]}, draw=black!55},
  vcap/.style={font=\tiny, text=black!60, align=center},
]
  \node[vcap,anchor=east] at (-0.42,1.55) {$\ff\in\bits^{|\ff|}$};
  \node (fflat) at (0.02,1.55)
    {$\begin{bmatrix}f_0\\f_1\\[-2pt]\vdots\\[-2pt]f_{|\ff|-1}\end{bmatrix}$};
  \node[vcap,anchor=east] at (-0.42,-0.15) {$\hh\in\bits^{|\hh|}$};
  \node (hflat) at (0.02,-0.15)
    {$\begin{bmatrix}h_0\\h_1\\[-2pt]\vdots\\[-2pt]h_{|\hh|-1}\end{bmatrix}$};
  \draw[arr] (fflat.east) -- node[vcap,above]{reshape} (1.65,1.55);
  \draw[arr] (hflat.east) -- node[vcap,above]{reshape} (1.65,-0.15);

  \begin{scope}[shift={(2.20,1.95)}]
    \foreach \r in {0,...,2}{
      \foreach \c in {0,...,4}{
        \node[cell,fill=blue!8] at (0.50*\c,-0.40*\r) {};
      }
    }
    \node[vcap] at (1.0,0.42) {$F$};
  \end{scope}

  \begin{scope}[shift={(2.20,0.25)}]
    \foreach \r in {0,...,2}{
      \foreach \c in {0,...,4}{
        \ifnum\c=2
          \node[cell,fill=orange!35] at (0.50*\c,-0.40*\r) {$h$};
        \else
          \node[cell,fill=orange!8] at (0.50*\c,-0.40*\r) {};
        \fi
      }
    }
    \node[vcap] at (1.0,0.42) {$H$};
    \draw[orange!75!black,rounded corners=1pt,line width=0.8pt]
      (0.73,0.22) rectangle (1.27,-1.02);
  \end{scope}

  \draw[arr,orange!75!black] (3.20,-0.85) -- (3.20,-1.35);
  \node[draw=orange!65!black,rounded corners=2pt,fill=orange!10,
        minimum height=0.58cm] (mu) at (3.20,-1.70)
    {$\mu[\colidx]=\ipZ{\gam}{H[*,\colidx]}$};
  \node[vcap,anchor=west,text width=3.25cm] at (5.15,-1.70)
    {fix one column; sum over $\rowidx$;\\the inner product is over $\Z$};
  \node[vcap,anchor=west,text width=3.25cm] at (5.15,1.55)
    {sumcheck weights:\\row $\eq(\cdot,r_{\rm row})\mapsto\gam$;\\
     column $c\,\eq(\cdot,r_{\rm col})$};
  \node[vcap] at (3.20,-2.35)
    {repeat independently for all $\ellcol$ columns};
\end{tikzpicture}
\captionof{figure}{The committed vector $\ff$ and virtual vector $\hh$ are
reshaped in their respective layouts into the data matrices $F$ and $H$.
The original public vector $\vv$ is not being factorised here:
\cref{step:eqshape} has already replaced it by the rank-one equality weights
$\eq(\cdot,r_{\rm row})$ and $c\,\eq(\cdot,r_{\rm col})$.
Each $\mu[\colidx]$ is a column fold: the selected column is fixed while the
sum runs over the row axis.}\label{fig:columnfold}
\end{center}

\begin{warning}[The integer messages need an explicit range check]
\label{warn:mujrange}
The field check \eqref{eq:murecombine} fixes $\mu[\colidx]$ only modulo $q$, while the next
step fixes $\gen^{\mu[\colidx]}$ only modulo $\ord(\gen)$.  Therefore the
verifier must reject unless
\[
  0\leq\mu[\colidx]\leq\ellrow(q-1),
  \qquad
  \ellrow(q-1)<\ord(\gen).
\]
The implementation performs this range check and verifies that $\gen$ is
primitive.  Without it, CRT lets a prover alter the field value while
preserving the exponent claim.
\end{warning}

\begin{warning}[Why the evaluation-shaping sumcheck is load-bearing]
\label{warn:rank1}
One cannot skip \cref{step:eqshape} and simply declare the reshaped matrix
$V=\operatorname{reshape}(\vv)$ to be rank one.  In general
$\vv=\bar A^{\!\top}\eq(\cdot,r_{\rm con})$ has no such factorisation; for the
shipped fold/clear splits its measured ranks on a schedule-shaped $A_0$ are
$1,2,4,4,3,2$ as $s-g$ ranges from $0$ to $5$.

The sumcheck fixes this for every axis split: its output coefficient is
$c\,\eq(\cdot,r_h)$, and equality factors exactly over
$r_h=(r_{\rm row},r_{\rm col})$.  A different, more expensive repair would be
to decompose the original table as
\[
  V[\rowidx,\colidx]
  =\sum_{a=1}^{r}v_{\mathrm{row}}^{(a)}[\rowidx]
        v_{\mathrm{col}}^{(a)}[\colidx],
\]
then run the integer-fold, forest, and GKR path for $r$ weight families.
Treating $r=1$ without either repair is a completeness error.
\end{warning}

\end{step}

% =========================================================================
\begin{step}{Move each integer column fold into an exponent}{%
  \clgdone{T}%
  \clgnew{P}{$\prod_{\rowidx}L_d[\rowidx,\colidx]=\gen^{\mu[\colidx]}$}%
  \clgnote{P}{one product per column, all in $\comfield$}%
}\label{step:exponent}
Fix a large-order generator $\gen\in\comfield^{\times}$ and define one leaf for
each matrix entry:
\[
  L_{\drow}[\rowidx,\colidx]
  \defeq
  1+\bigl(\gen^{\gam[\rowidx]}-1\bigr)
       \pitwo(H[\rowidx,\colidx]).                \tag{4.5}\label{eq:leafdef}
\]
Because $H[\rowidx,\colidx]$ is a bit, the leaf is $1$ when the bit is zero and
$\gen^{\gam[\rowidx]}$ when it is one.  Therefore, for every fixed column,
\[
  \boxed{
  \prod_{\rowidx\in\Rows}L_{\drow}[\rowidx,\colidx]
  =\gen^{\sum_{\rowidx}\gam[\rowidx]H[\rowidx,\colidx]}
  =\gen^{\mu[\colidx]}.}                         \tag{4.6}
\]
The product iterates over $\rowidx$; $\colidx$ selects which of the
$\ellcol$ products is being proved.  This is not a sumcheck.  It is the
algebraic bridge from addition over $\Z$ to multiplication in
$\comfield^{\times}$.

\end{step}

% =========================================================================
\begin{step}{Reduce all column products with one merged GKR forest}{%
  \clgdone{P}%
  \clgnew{L}{$e_{\drow}+1=\sum_{\rowidx,\colidx}\!R\,\pitwo(H)$}%
  \clgnote{L}{one leaf-boundary claim over $\comfield$}%
}\label{step:gkr}
There is one product tree for every column, but the trees are proved as a
single forest.  At level $k$, where the root level is $k=0$ and the leaves are
$k=\drow$, use
\[
  \rowprefix\in\bits^k,\qquad \colidx\in\Cols,
\]
and define the recursion
\[
  \boxed{
  L_k[\rowprefix,\colidx]
  =L_{k+1}[\rowprefix\Vert0,\colidx]\,
   L_{k+1}[\rowprefix\Vert1,\colidx].}            \tag{4.7}
\]
The root table is public from Step 5's integers:
\[
  L_0[\colidx]=\gen^{\mu[\colidx]}.
\]
The verifier samples $\zcol^{(0)}\in\comfield^{\dcol}$ and starts from
\[
  E_0^{\mathrm{GKR}}:\quad
  \widetilde L_0(\zcol^{(0)})=e_0,
  \qquad
  e_0=\sum_{\colidx\in\Cols}
       \eq(\colidx,\zcol^{(0)})\gen^{\mu[\colidx]}.
\]

\paragraph{The claim proved at every level.}
Suppose level $k$ starts with
\[
  E_k^{\mathrm{GKR}}:\quad
  \widetilde L_k(\zrow^{(k)},\zcol^{(k)})=e_k,
  \qquad \zrow^{(k)}\in\comfield^k.
\]
The layer sumcheck proves
\begin{align}
  e_k
  ={}&\sum_{\rowprefix\in\bits^k}
      \sum_{\colidx\in\Cols}
      \eq(\rowprefix,\zrow^{(k)})
      \eq(\colidx,\zcol^{(k)})                         \notag\\[-2pt]
     &\hspace{2.4cm}\cdot
      L_{k+1}[\rowprefix\Vert0,\colidx]
      L_{k+1}[\rowprefix\Vert1,\colidx].              \tag{4.8}\label{eq:gkrlayer}
\end{align}
The implementation realizes \eqref{eq:gkrlayer} as two consecutive
degree-$3$ sumchecks.

\paragraph{Which variables are sumchecked at depth $k$.}
Write a complete leaf-row index as
$x=(x_1,\ldots,x_{\drow})$.  A node at depth $k$ retains only the
already-exposed prefix $x_{\leq k}=(x_1,\ldots,x_k)$.  Therefore the first of
the two sumchecks, the row-prefix phase, uses exactly the iteration variables
\[
  X=(X_1,\ldots,X_k)\in\bits^k,
\]
not the unexposed suffix $(x_{k+1},\ldots,x_{\drow})$.  It consequently has
$k$ rounds: none at the root, one at depth $1$, and one additional round at
each deeper multiplication layer.  The child arguments $X\Vert0$ and
$X\Vert1$ append the next row coordinate.

For this row-prefix sumcheck, define the
column-contracted degree-$2$ polynomial
\[
  G_k(X)
  \defeq
  \sum_{\colidx\in\Cols}
    \eq(\colidx,\zcol^{(k)})
    \mle{L_{k+1}}(X,0,\colidx)
    \mle{L_{k+1}}(X,1,\colidx).
\]
The row-prefix sumcheck starts from
\[
  e_k=\sum_{\rowprefix\in\bits^k}
       \eq(\rowprefix,\zrow^{(k)})G_k(\rowprefix).
                                                        \tag{4.8a}
\]
After its row-prefix challenges fix
$r_{\mathrm{row}}^{(k)}$, let
$s_k=G_k(r_{\mathrm{row}}^{(k)})$.  If
$e_{k,\mathrm{row}}^{\mathrm{end}}$ is this sumcheck's terminal value, the
verifier checks
\[
  e_{k,\mathrm{row}}^{\mathrm{end}}
  =\eq(r_{\mathrm{row}}^{(k)},\zrow^{(k)})s_k,
\]
and the column-index sumcheck starts from
\[
  s_k=\sum_{\colidx\in\Cols}
       \eq(\colidx,\zcol^{(k)})
       E_k[\colidx]O_k[\colidx],                       \tag{4.8b}
\]
where
\[
  E_k[\colidx]
   =\widetilde L_{k+1}(r_{\mathrm{row}}^{(k)},0,\colidx),
  \qquad
  O_k[\colidx]
   =\widetilde L_{k+1}(r_{\mathrm{row}}^{(k)},1,\colidx).
\]
Thus the two groups of iteration variables are:
\[
  \underbrace{k\text{ coordinates of }\rowprefix}_{\text{row-prefix sumcheck}}
  \qquad\text{then}\qquad
  \underbrace{\dcol\text{ coordinates of }\colidx}_{\text{column-index sumcheck}}.
\]
The row-prefix sumcheck is empty at the root.  The column-index sumcheck is what
merges all columns; without it, the prover would expose one child pair per
column at every level.

At the end of the level, the verifier has child values
\[
  a_k=\widetilde L_{k+1}(r_{\mathrm{row}}^{(k)},0,
                         r_{\mathrm{col}}^{(k)}),
  \qquad
  b_k=\widetilde L_{k+1}(r_{\mathrm{row}}^{(k)},1,
                         r_{\mathrm{col}}^{(k)}).
\]
If $e_{k,\mathrm{col}}^{\mathrm{end}}$ is the column-index sumcheck's terminal
value, the verifier checks
\[
  e_{k,\mathrm{col}}^{\mathrm{end}}
  =\eq(r_{\mathrm{col}}^{(k)},\zcol^{(k)})a_kb_k.
\]
Only then does the transcript absorb $(a_k,b_k)$ and sample
$\chi_k\in\comfield$, folding the pair into the next boundary claim:
\[
  \boxed{
  \begin{aligned}
  \zrow^{(k+1)}&=(r_{\mathrm{row}}^{(k)},\chi_k),\\
  \zcol^{(k+1)}&=r_{\mathrm{col}}^{(k)},\\
  e_{k+1}&=(1-\chi_k)a_k+\chi_k b_k,\\
  E_{k+1}^{\mathrm{GKR}}&:
    \widetilde L_{k+1}(\zrow^{(k+1)},\zcol^{(k+1)})=e_{k+1}.
  \end{aligned}}                                      \tag{4.9}
\]
This is the recursive relation: each level consumes one claim about $L_k$ and
creates one claim about $L_{k+1}$.

\Cref{fig:gkr-depth4} shows the complete descent for a depth-$4$ example.  The
tree is representative: the row-prefix sumcheck moves within all $2^s$ trees,
whereas the column-index sumcheck is the horizontal contraction that merges
their column index.  In \eqref{eq:gkrlayer}, $\rowprefix\Vert b$ appends the
newly exposed top coordinate.  The figure uses recursive tree order for
legibility, not the low-bit-first flat-table order used by the implementation.

\begin{center}
\begin{tikzpicture}[
  x=1cm,
  y=1cm,
  font=\scriptsize,
  levelband/.style={fill=blue!3, rounded corners=1pt},
  treeedge/.style={draw=black!25, line width=0.35pt},
  treenode/.style={circle, draw=black!42, fill=black!4,
                   minimum size=2.8mm, inner sep=0pt, line width=0.35pt},
  leafnode/.style={treenode, minimum size=1.8mm, fill=orange!9},
  claim/.style={draw=blue!55!black, fill=blue!4, rounded corners=1.5pt,
                minimum width=7.15cm, minimum height=0.42cm,
                inner sep=1.1pt, line width=0.45pt},
  rowsc/.style={draw=blue!55!black, fill=blue!10, rounded corners=1.2pt,
                minimum width=2.20cm, minimum height=0.48cm,
                align=center, inner sep=1pt, line width=0.35pt},
  colsc/.style={draw=orange!70!black, fill=orange!12, rounded corners=1.2pt,
                minimum width=2.45cm, minimum height=0.48cm,
                align=center, inner sep=1pt, line width=0.35pt},
  linefold/.style={draw=black!45, fill=black!4, rounded corners=1.2pt,
                   minimum width=1.55cm, minimum height=0.48cm,
                   align=center, inner sep=1pt, line width=0.35pt},
  claimarr/.style={-{Latex[length=1.5mm]}, draw=blue!50!black,
                   line width=0.45pt},
  tinyarr/.style={-{Latex[length=1.2mm]}, draw=black!45,
                   line width=0.35pt},
  rowdetail/.style={draw=blue!55!black, fill=blue!7, rounded corners=1.5pt,
                    text width=3.80cm, minimum height=1.68cm,
                    align=center, inner sep=3pt, line width=0.4pt},
  coldetail/.style={draw=orange!70!black, fill=orange!8,
                    rounded corners=1.5pt, text width=3.75cm,
                    minimum height=1.68cm, align=center,
                    inner sep=3pt, line width=0.4pt},
  folddetail/.style={draw=black!45, fill=black!3, rounded corners=1.5pt,
                     text width=3.35cm, minimum height=1.68cm,
                     align=center, inner sep=3pt, line width=0.4pt},
  boundary/.style={draw=orange!70!black, fill=orange!7,
                   rounded corners=1.5pt, align=center,
                   inner sep=3pt, line width=0.45pt},
  dcap/.style={font=\tiny, text=black!58, align=center},
]

% One continuous swimlane: every level table is horizontally aligned with the
% verifier claim about that table.  Operations in the gap reduce that claim to
% the claim on the next level.
\node[dcap, text width=4.1cm] at (2.13,0.70)
  {one representative column tree; the whole tree is repeated $2^s$ times};
\node[dcap, text width=7.15cm] at (8.28,0.70)
  {each level claim is followed by its row-prefix sumcheck, column-index sumcheck,
   and line fold, side by side};

\draw[draw=black!22, dashed, rounded corners=2pt, line width=0.35pt]
  (0.10,0.30) rectangle (4.22,-4.62);
\node[dcap, fill=white, inner sep=0.7pt] at (2.16,-4.62)
  {one complete product tree};

\foreach \lev in {0,...,4} {
  \pgfmathtruncatemacro{\numnodes}{2^\lev}
  \pgfmathtruncatemacro{\lastnode}{\numnodes-1}
  \pgfmathsetmacro{\yy}{-1.10*\lev}
  \path[levelband] (0.18,\yy-0.17) rectangle (4.14,\yy+0.17);
  \foreach \idx in {0,...,\lastnode} {
    \pgfmathsetmacro{\xx}{0.24+3.84*(\idx+0.5)/\numnodes}
    \coordinate (gkr-\lev-\idx) at (\xx,\yy);
  }
  \ifnum\lev=4
    \node[dcap, anchor=east] at (0.06,\yy) {$L_4=\Lambda$};
  \else
    \node[dcap, anchor=east] at (0.06,\yy) {$L_{\lev}$};
  \fi
}

% Recursive visual order keeps the binary tree legible.  The surrounding prose
% states that the implementation's flat tables instead use low-bit-first order.
\foreach \lev in {0,...,3} {
  \pgfmathtruncatemacro{\numnodes}{2^\lev}
  \pgfmathtruncatemacro{\lastnode}{\numnodes-1}
  \pgfmathtruncatemacro{\childlev}{\lev+1}
  \foreach \idx in {0,...,\lastnode} {
    \pgfmathtruncatemacro{\leftchild}{2*\idx}
    \pgfmathtruncatemacro{\rightchild}{2*\idx+1}
    \draw[treeedge] (gkr-\lev-\idx) -- (gkr-\childlev-\leftchild);
    \draw[treeedge] (gkr-\lev-\idx) -- (gkr-\childlev-\rightchild);
  }
}

\foreach \lev in {0,...,4} {
  \pgfmathtruncatemacro{\numnodes}{2^\lev}
  \pgfmathtruncatemacro{\lastnode}{\numnodes-1}
  \foreach \idx in {0,...,\lastnode} {
    \ifnum\lev=4
      \node[leafnode] at (gkr-\lev-\idx) {};
    \else
      \node[treenode] at (gkr-\lev-\idx) {};
    \fi
  }
}

\draw[-{Latex[length=1.7mm]}, draw=black!38, line width=0.45pt]
  (-0.36,-4.22) -- (-0.36,-0.18);
\node[dcap, rotate=90, anchor=center] at (-0.68,-2.20)
  {products computed upward};

% Claims live at the same heights as the product-tree levels.
\node[claim] (E0) at (8.28, 0.00)
  {$E_0^{\mathrm{GKR}}:\ \mle{L_0}(\zcol^{(0)})=e_0$};
\node[claim] (E1) at (8.28,-1.10)
  {$E_1^{\mathrm{GKR}}:\ \mle{L_1}(\zrow^{(1)},\zcol^{(1)})=e_1$};
\node[claim] (E2) at (8.28,-2.20)
  {$E_2^{\mathrm{GKR}}:\ \mle{L_2}(\zrow^{(2)},\zcol^{(2)})=e_2$};
\node[claim] (E3) at (8.28,-3.30)
  {$E_3^{\mathrm{GKR}}:\ \mle{L_3}(\zrow^{(3)},\zcol^{(3)})=e_3$};
\node[claim] (E4) at (8.28,-4.40)
  {$E_4^{\mathrm{GKR}}:\ \mle{L_4}(\zrow^{(4)},\zcol^{(4)})=e_4$};

\foreach \k/\next/\rowrounds/\midy in {
  0/1/0/-0.55,
  1/2/1/-1.65,
  2/3/2/-2.75,
  3/4/3/-3.85} {
  \node[rowsc] (R\k) at (5.70,\midy)
    {row-prefix sumcheck\\$\rowrounds$ rounds};
  \node[colsc] (C\k) at (8.28,\midy)
    {column-index sumcheck\\$s$ rounds};
  \node[linefold] (X\k) at (10.82,\midy)
    {$\chi_{\k}$ line fold};
  \draw[claimarr] (E\k.south west) -- (R\k.north west);
  \draw[tinyarr] (R\k) -- (C\k);
  \draw[tinyarr] (C\k) -- (X\k);
  \draw[claimarr] (X\k.south east) -- (E\next.north east);
}

% Light guides make the level-table / claim correspondence explicit without
% implying that a single selected tree path is being verified.
\foreach \lev in {0,...,4} {
  \draw[draw=black!20, densely dotted, line width=0.3pt]
    (4.22,-1.10*\lev) -- (E\lev.west);
}

% Expand one transition into its exact three consecutive relations.  The boxes
% are horizontal because these are stages of one claim reduction, not functions
% named A and B.
\node[rowdetail] (rowdetail) at (2.05,-6.30) {%
  \textbf{Row-prefix sumcheck} \ ($k$ rounds)\\[2pt]
  $\begin{aligned}
    e_k&=\displaystyle\sum_{x\in\bits^k}
      \eq(x,\zrow^{(k)})\\[-1pt]
      &\qquad\cdot G_k(x),\\[2pt]
    e_{k,\mathrm{row}}^{\mathrm{end}}
      &=\eq(r_{\mathrm{row}}^{(k)},\zrow^{(k)})s_k
  \end{aligned}$
};

\node[coldetail] (coldetail) at (6.25,-6.30) {%
  \textbf{Column-index sumcheck} \ ($s$ rounds)\\[2pt]
  $\begin{aligned}
    s_k&=\displaystyle\sum_{c\in\bits^s}
      \eq(c,\zcol^{(k)})\\[-1pt]
      &\qquad\cdot E_k(c)O_k(c),\\[2pt]
    e_{k,\mathrm{col}}^{\mathrm{end}}
      &=\eq(r_{\mathrm{col}}^{(k)},\zcol^{(k)})a_kb_k
  \end{aligned}$
};

\node[folddetail] (folddetail) at (10.34,-6.30) {%
  \textbf{Absorb the child pair; then fold}\\[2pt]
  $\begin{gathered}
    e_{k+1}=a_k+\chi_k(a_k+b_k),\\[2pt]
    \zrow^{(k+1)}=(r_{\mathrm{row}}^{(k)},\chi_k),\\[1pt]
    \zcol^{(k+1)}=r_{\mathrm{col}}^{(k)}
  \end{gathered}$
};

\draw[claimarr] (rowdetail) -- (coldetail);
\draw[claimarr] (coldetail) -- (folddetail);

\node[boundary, text width=12.05cm] at (6.14,-8.62) {%
  $\displaystyle
  \begin{aligned}
    e_4+1
      &=\sum_{i\in\bits^4}\sum_{c\in\bits^s}
        \eq(i,\zrow^{(4)})\eq(c,\zcol^{(4)})
        \bigl(\gen^{\gam[i]}-1\bigr)\pitwo(H[i,c])\\[-1pt]
      &\hspace{5.3cm}\longrightarrow\quad
        \text{row-only pre-sumcheck}.
  \end{aligned}$
};
\end{tikzpicture}
\captionof{figure}{A depth-$4$ merged GKR forest drawn as one claim-reduction
swimlane.  Each level table $L_k$ is horizontally aligned with its open claim
$E_k^{\mathrm{GKR}}$.  Between levels, the row-prefix and column-index
sumchecks are displayed side by side, followed by the line fold that creates
the next claim.  The row-prefix sumcheck has $k$ degree-$3$ rounds at transition
$k$ and is empty at the root; the column-index sumcheck always has $s$ rounds.
The displayed tree uses recursive visual order rather than the implementation's
low-bit-first flat-table order.  Across all four transitions there are $6+4s$
degree-$3$ sumcheck rounds and four line folds.}
\label{fig:gkr-depth4}
\end{center}


\paragraph{The final leaf boundary.}
After $\drow$ levels, write the final points simply as $\zrow$ and $\zcol$.
Substituting the affine leaf \eqref{eq:leafdef} gives
\begin{align*}
  e_{\drow}
    &=1+
      \sum_{\rowidx\in\Rows}\sum_{\colidx\in\Cols}
      \eq(\rowidx,\zrow)\eq(\colidx,\zcol)
      \bigl(\gen^{\gam[\rowidx]}-1\bigr)
      \pitwo(H[\rowidx,\colidx]).
\end{align*}
Since $\comfield$ has characteristic two, moving the constant to the other side
also writes it as $+1$:
\[
  \boxed{\begin{aligned}
  L:\quad e_{\drow}+1
  ={}&\sum_{\rowidx\in\Rows}\sum_{\colidx\in\Cols}
      \eq(\rowidx,\zrow)\eq(\colidx,\zcol)\\[-2pt]
     &\quad\cdot\bigl(\gen^{\gam[\rowidx]}-1\bigr)
      \pitwo(H[\rowidx,\colidx]).
  \end{aligned}}                                  \tag{4.10}\label{eq:leafboundary}
\]
The two displayed sums in \eqref{eq:leafboundary} are finite matrix-axis sums:
$\ellrow$ row values and $\ellcol$ column values.  The sumcheck rounds in
\eqref{eq:gkrlayer}, by contrast, iterate over the binary coordinates that encode those
indices.

\begin{remark}[Why the leaf weight is shared across columns]
\label{rem:sharedv}
The non-$\eq$ factor $\gen^{\gam[\rowidx]}-1$ depends only on the row.  The
merged descent also produces one shared row point $\zrow$.  Consequently every
column reaches the boundary with the same row-weight table; only
$\eq(\colidx,\zcol)$ changes across columns.
\end{remark}

\end{step}

% =========================================================================
\begin{step}{Run the row-only pre-sumcheck}{%
  \clgdone{L}%
  \clgnew{$E_H$}{$\widetilde H(r_{\rm row}^{*},\zcol)=\eta$}%
  \clgnote{$E_H$}{$\drow$ sumcheck rounds; no column variable remains}%
}\label{step:presum}
First contract the already-fixed column point:
\[
  H_{\zcol}[\rowidx]
  \defeq
  \sum_{\colidx\in\Cols}
    \eq(\colidx,\zcol)\pitwo(H[\rowidx,\colidx]),
\]
and define the public row table
\[
  R[\rowidx]
  \defeq
  \eq(\rowidx,\zrow)
  \bigl(\gen^{\gam[\rowidx]}-1\bigr).
\]
Then the GKR boundary is exactly
\[
  e_{\drow}+1
  =\ipC{R}{H_{\zcol}}
  =\sum_{\rowidx\in\Rows}R[\rowidx]H_{\zcol}[\rowidx].  \tag{4.11}\label{eq:presumstart}
\]
Run a degree-$2$ sumcheck on \eqref{eq:presumstart}.  Its iteration variables are the
$\drow$ Boolean coordinates encoding $\rowidx$:
\[
  \rho_1,\ldots,\rho_{\drow}.
\]
There are \textbf{no column-index sumcheck variables}: GKR has already fixed
$\zcol$, and the column contraction is part of the multilinear target.  The
sumcheck ends at a fresh point
$r_{\mathrm{row}}^{*}\in\comfield^{\drow}$ and leaves
\[
  \boxed{
  E_H:\quad
  \widetilde H(r_{\mathrm{row}}^{*},\zcol)=\eta.}      \tag{4.12}
\]
The verifier evaluates the known factor $\widetilde R(r_{\mathrm{row}}^{*})$
in the final sumcheck check.

\begin{remark}[The two sumchecks do different jobs]
\label{rem:deferredeq}
\Cref{step:eqshape} runs over $\Fq$ before the integer folds.  It converts the
arbitrary PIOP coefficient $\vv$ into the separable weight
$c\,\eq(\cdot,r_h)$; without it, Step 5 is not defined in general.  This
row-only sumcheck runs later over $\comfield$, after GKR, and converts the
nonlinear leaf boundary into one evaluation of $H$ so that the virtual map can
be transposed.  Neither sumcheck subsumes the other.
\end{remark}

\end{step}

% =========================================================================
\begin{step}{Transpose the full virtual map and eliminate $\hh$}{%
  \clgdone{$E_H$}%
  \clgdone{H}%
  \clgnew{$E_f$}{$\sum_a\lambda_a\widetilde F(r_f^{(a)})=\eta+\kappa_0$}%
  \clgnote{$E_f$}{only committed bits remain}%
}\label{step:transpose}
Combine the Step 8 point as
\[
  r_H=(r_{\mathrm{row}}^{*},\zcol),
  \qquad
  q_H[\mathsf{pos}]=\eq(\mathsf{pos},r_H).
\]
The evaluation claim is the field-annotated inner product
\[
  \eta=\ipC{q_H}{\pitwo(\hh)}.
\]
Now use the standing relation $H$ in its native characteristic:
\[
  \boxed{
  \eta
   =\ipC{q_H}{M\pitwo(\ff^{\circ})}
   =\ipC{M^{\!\top}q_H}{\pitwo(\ff^{\circ})}.}        \tag{4.13}\label{eq:fulladjoint}
\]
This identity is valid in $\comfield$ because it contains $\Ftwo$ and
$\pitwo(\hh)=M\pitwo(\ff^{\circ})$ there.  It would be false over $\Z$ or
$\Fq$, where lifting an XOR is nonlinear.

Split the adjoint vector according to $\ff^{\circ}=(1,\ff)$:
\[
  M^{\!\top}q_H=(\kappa_0,\uu),
  \qquad
  \ipC{\uu}{\pitwo(\ff)}=\eta+\kappa_0.               \tag{4.14}\label{eq:splitconstant}
\]
The scalar $\kappa_0$ is the known contribution of the constant coordinate.
This is where the public/affine part is absorbed; it is not a separate public
fold in the PIOP.

For SHA-256, the full transpose is represented by the tap closure rather than
materializing $\uu$.  It expresses the adjoint as a short structured sum of
evaluation vectors,
\[
  \uu=\sum_a\lambda_a\eq(\cdot,r_f^{(a)}),
\]
so with $F$ the MLE of the committed bit table, \eqref{eq:splitconstant} becomes
\[
  \boxed{
  E_f:\quad
  \sum_a\lambda_a\widetilde F(r_f^{(a)})
  =\eta+\kappa_0.}                                  \tag{4.15}\label{eq:committedclaim}
\]
No virtual coordinate remains.

\begin{warning}[Use the full SHA-256 adjoint, not a columnwise $M_0$]
\label{warn:mblockdiag}
The tempting shortcut
$M=\mathrm{Id}_{\mathrm{col}}\otimes M_0$ is false for SHA-256.  Rotations
change the clear bit-position coordinate, while schedule and state taps can
carry between the row and column parts of the flattened index.  Therefore a
formula that transposes one identical $M_0$ independently inside every column
does not discharge $H$.  Equation \eqref{eq:fulladjoint} is the always-correct statement;
the implementation realizes its clear-index-dependent support with
\src{tap\_closure\_desc} and \src{residual\_b\_evals\_tap --- src/taps.rs}.
Any optimization must be proved equal to the full $M^{\!\top}q_H$.
\end{warning}

\end{step}

% =========================================================================
\begin{step}{Ring-switch from bit evaluations to packed evaluations}{%
  \clgdone{$E_f$}%
  \clgnew{$E_P$}{packed MLE evaluations}%
  \clgnote{$E_P$}{$128=2^7$ bit coordinates per $\comfield$ symbol}%
}\label{step:ringswitch}
The commitment encodes
\[
  P=\pack(\pitwo(\ff))\in\comfield^{|\ff|/128},
\]
not the raw bit table $F$.  For each evaluation in $E_f$, split the point so
the low seven row-axis coordinates select one of the $128$ bits inside a packed
symbol.  The Galois ring-switch expands the bit evaluation into $2^7$ partial
evaluations of the packed MLE and recombines them with an $\Ftwo$-linear map.
Applying that map to all terms in \eqref{eq:committedclaim} yields a finite batch
\[
  \boxed{
  E_P:\quad
  \sum_b\alpha_b\widetilde P(r_P^{(b)})=\kappa}       \tag{4.16}
\]
whose coefficients and target are public transcript values.

Packing is along the row axis and may never cross a column boundary.  This is
why $128\mid\ellrow$ and why the packing width is structural rather than a
tunable serialization choice.

\end{step}

% =========================================================================
\begin{step}{Open the packed evaluations against the oracle}{%
  \clgdone{$E_P$}%
  \clgdone{C}%
  \clgnote{C}{recursive Ligerito proves proximity and the evaluations}%
}\label{step:proximity}
Recursive Ligerito proves that the oracle is close to an encoding of one packed
message $P$ and that the requested MLE evaluations satisfy $E_P$.  The same
opening therefore closes two obligations at once:
\[
  \orc=\Enc_{\cC}(P)
  \qquad\text{and}\qquad
  \sum_b\alpha_b\widetilde P(r_P^{(b)})=\kappa.
\]
This is the only step that queries the Merkle-authenticated oracle.  Its own
recursive degree-$2$ sumchecks and proximity queries supply the remaining
soundness error.

\end{step}

\noindent
The two phases have a clean boundary.  Steps 1--4 use the public constraint
matrix and one degree-$2$ sumcheck to turn $A\hh=0$ into an evaluation-shaped
field claim.  Steps 5--11 convert that claim to integer column folds, prove
their products, shape the GKR leaf into an
evaluation, eliminate the virtual witness with the full transpose, ring-switch
to the committed packing, and finally open the oracle.

\subsection{The complete eleven-step chain}

\begin{center}\small
\begin{tabular}{@{}clp{0.43\textwidth}l@{}}
\toprule
step & phase & claim after the step & ambient algebra\\
\midrule
1 & PIOP & $Q(\beta)=0$ on $\bits^{\nu}$ & $\Z$\\
2 & PIOP & $\bar A\pi_q(\hh)=0$ & $\Fq$\\
3 & PIOP & $E_0:\ipq{\vv}{\pi_q(\hh)}=0$ & $\Fq$\\
4 & PIOP & $E_{\rm eq}:\ipq{c\,\eq(\cdot,r_h)}{\pi_q(\hh)}=\tau$ & $\Fq$\\
5 & IOP & $T:\mu[\colidx]=\ipZ{\gam}{H[*,\colidx]}$ & $\Z$\\
6 & IOP & $P:\prod_{\rowidx}L_{\drow}=\gen^{\mu[\colidx]}$ & $\comfield^\times$\\
7 & IOP & $L:$ the leaf-boundary relation \eqref{eq:leafboundary} & $\comfield$\\
8 & IOP & $E_H:\widetilde H(r_{\rm row}^{*},\zcol)=\eta$ & $\comfield$\\
9 & IOP & $E_f:\sum_a\lambda_a\widetilde F(r_f^{(a)})=\eta+\kappa_0$ & $\comfield$\\
10 & IOP & $E_P:$ packed-MLE evaluations & $\comfield$\\
11 & IOP & no open claims & $\comfield$\\
\bottomrule
\end{tabular}
\end{center}

\subsection{Matrix-shape trade-offs}
\label{sec:matrix-shape-tradeoffs}

Hold the padded single-bit table size fixed:
\[
  \ell=\ellrow\ellcol,
  \qquad
  \ellrow=2^{\drow},
  \qquad
  \ellcol=2^{\dcol}.
\]
For SHA-256, moving one Boolean coordinate from the row axis to the column axis
changes the shape by
\[
  (\ellrow,\ellcol,\drow,\dcol)
  \longmapsto
  (\ellrow/2,\,2\ellcol,\,\drow-1,\,\dcol+1).       \tag{4.17}
\]
Thus the forest has twice as many trees, each half as tall.  This is a reshape,
not a compression: the prover still touches essentially all $\ell$ entries.
Ignoring live-zero-column elision, define the merged-GKR round count
\[
  Q(\drow,\dcol)
  \defeq
  \sum_{k=0}^{\drow-1}(k+\dcol)
  =\frac{\drow(\drow-1)}2+\drow\dcol.                \tag{4.18}
\]
Assume an admissible arity-$2$ move with $128\mid\ellrow/2$.
\Cref{step:eqshape} guarantees rank-one equality weights for every such split.
Fold/forest/GKR/pre-sumcheck rows are per chunk;
commitment/Ligerito are shared, and memory is peak live storage.

\begingroup
\footnotesize
\renewcommand{\arraystretch}{1.10}
\begin{longtable}{@{}
  >{\raggedright\arraybackslash}p{0.145\textwidth}
  >{\raggedright\arraybackslash}p{0.215\textwidth}
  >{\raggedright\arraybackslash}p{0.235\textwidth}
  >{\raggedright\arraybackslash}p{0.225\textwidth}@{}}
\caption{Shape trade-offs at fixed padded bit count.  More columns create more,
shallower trees but do not remove the prover's linear pass.}
\label{tab:matrix-shape-tradeoffs}\\
\toprule
component & prover at fixed $\ell$ & verifier or proof & when $\ellcol$ doubles \\
\midrule
\endfirsthead
\multicolumn{4}{@{}l}{\footnotesize\emph{Table
\thetable\ continued from the previous page.}}\\
\toprule
component & prover at fixed $\ell$ & verifier or proof & when $\ellcol$ doubles \\
\midrule
\endhead
\midrule
\multicolumn{4}{r@{}}{\footnotesize\emph{continued on the next page}}\\
\endfoot
\bottomrule
\endlastfoot
evaluation-shaping sumcheck
  & one $\Theta(\ell)$ pass over the padded $h$ table
  & $\drow+\dcol=d_h$ degree-$2$ rounds, plus the public evaluation $c=\mle{\vv}(r_h)$
  & unchanged: the total dimension and padded table length are fixed \\

integer column folds
  & $\ellcol$ folds of length $\ellrow$: $\Theta(\ell)$ total work
  & $\Theta(\ellcol)$ checks and recombination; $16\ellcol$ bytes per chunk
  & the prover scan is unchanged; the sent fold block doubles \\

product forest
  & $\ellcol(\ellrow-1)=\ell-\ellcol$ internal products; depth $\drow$
  & derive $\ellcol$ roots $\gen^{\mu[\colidx]}$ and evaluate their MLE: $\Theta(\ellcol)$
  & one level disappears and exactly the old $\ellcol$ internal products disappear; root work doubles \\

merged GKR
  & row-prefix phases are $\Theta(\ell)$; column-index phases contribute $\Theta(\drow\ellcol)$
  & $Q$ degree-$3$ rounds; GKR messages are about $48Q+32\drow$ bytes
  & $Q$ falls by exactly $\dcol$; tree parallelism grows, but every column phase is twice as wide \\

row-only pre-sumcheck
  & the column contraction remains $\Theta(\ell)$; the residual row table has length $\ellrow$
  & a length-$\ellrow$ weight table and $\drow$ degree-$2$ rounds
  & the contraction is unchanged; the row table halves and one round disappears \\

commitment and Ligerito
  & one encoding pass over the fixed packed-message length
  & one commitment and one final Ligerito proof are shared across all $L$ chunks
  & unchanged at fixed $\ell$ \\

memory and cache
  & lazy-forest peak is $\Theta(\ell)$; shared row tables are $O(\ellrow)$ and column buffers $O(\ellcol)$
  & row-weight storage is $O(\ellrow)$; roots and sent folds are $O(\ellcol)$
  & leading peak is similar; shared row tables halve while column buffers double \\

exponent range and chunks
  & the unchunked bound $\ellrow(q-1)$ halves; at $W=1$, $c_w=126-\drow$ gains one bit and may lower $L$
  & each chunk remains $<2^{127}$; the forest-through-ring-switch part scales roughly with $L$
  & a large discontinuous win only if $L$ falls; the usual SHA-256 shape has $L=1$ \\
\end{longtable}
\endgroup

The trade is especially visible at $\ell=2^{28}$.  Changing
$(\drow,\dcol)=(17,11)$ to $(16,12)$ halves the row tables and changes
$Q=323$ to $Q=312$, but doubles the root count and the sent-fold block from
$32$ to $64$ KiB.  For this $r=L=1$ comparison, the leading shape-dependent
cost is
\[
  \boxed{
    \text{prover}=\Theta(\ell),
    \qquad
    \text{verifier}=\Theta(\ellrow+\ellcol),
    \qquad
    \text{sent folds}=\Theta(\ellcol).}              \tag{4.19}
\]
Consequently verifier work alone is minimised near
$\ellrow\approx\ellcol\approx\sqrt\ell$.  The dominant sent-fold term favours fewer columns,
whereas shallower trees can improve prover wall-clock time through smaller
shared tables, better cache behaviour and more independent trees.  The split is
therefore a balance, not a free prover speedup; \cref{sec:config} records the
SHA-256-specific layout pins and the recommended depth cap.
